\chapter{Introduction}
\label{ch:introduction}

When a computation is delegated to a party we do not fully trust, we
would like a way to check that it was performed correctly without
redoing the work ourselves. Proof systems make this possible: a prover
produces a short certificate that convinces a verifier of the
correctness of a statement, and the verifier checks it far more cheaply
than recomputing the statement from scratch. Over the last decade these
systems have moved from theory to practice, underpinning scalable and
privacy-preserving applications, and much of their machinery rests on a
single building block, the \emph{polynomial commitment scheme}.

These systems are no longer only a theoretical construction, and the
polynomial commitment scheme is the component on which a great many of
them are built. The most visible application is the scaling of
blockchains through \emph{zk-rollups}: instead of having every node
re-execute thousands of transactions, a single succinct proof attests
that the whole batch was processed correctly, and the base chain only
verifies that proof. StarkNet, a Layer-2 network on Ethereum, does
exactly this using STARK proofs, which are built on top of the FRI
protocol studied in this thesis; other systems such as zkSync and
Polygon zkEVM follow the same idea. A second family of applications
exploits the fact that a proof can reveal that a statement is true while
hiding the data behind it: this enables privacy-preserving payments,
identity and credential systems that prove a fact (for instance that a
person is over a certain age) without disclosing the underlying document,
and privacy-preserving analysis of sensitive data such as medical
records. A third use is \emph{verifiable computation}, where an
expensive calculation is delegated to an untrusted server that returns
the result together with a proof of its correctness.

FRI in particular has found uses beyond proving computations. A recent
example is \emph{data availability sampling}~\cite{FRIDA}, where a light
node of a blockchain, which does not store the block contents, wants to
be sure that those contents are actually available in the network
without downloading them in full. By probing only a few random positions
of an encoded block, and checking them against a short commitment, the
node becomes convinced that the whole block can be recovered. Hall-Andersen,
Simkin, and Wagner show that FRI can be turned into exactly such a
scheme, which illustrates how the same low-degree testing machinery
studied in this thesis serves as a building block far beyond its original
purpose. What all these applications share is the same core need: a way
to commit to data and later prove statements about it, cheaply and
reliably.

A polynomial commitment scheme lets a prover commit to a polynomial with
a short digest and later prove its value at any chosen point, without
revealing the polynomial itself. Different schemes make different
trade-offs. Some rely on a trusted setup or on elliptic-curve
assumptions that are broken by quantum computers. The scheme studied in
this thesis avoids both: it is built on FRI, the Fast Reed-Solomon
Interactive Oracle Proof of Proximity, and needs only a hash function.
It is therefore \emph{transparent}, requiring no trusted setup, and
plausibly \emph{post-quantum}, since its security rests on hashing alone.

\section{Contribution}
\label{sec:contribution}

This thesis builds such a scheme from the ground up, over the BabyBear
prime field $\mathbb{F}_p$ with $p = 2^{31} - 2^{27} + 1 = 15 \cdot 2^{27} + 1 = 2013265921$,
and implements every part of it from scratch in Python
without external cryptographic libraries. This includes the field
arithmetic and its quartic extension, the Number Theoretic Transform,
the Merkle commitments, a Keccak-based duplex-sponge challenger, and the
FRI prover and verifier, on top of which the polynomial commitment
scheme is realised through the quotient technique. The complete source code
is publicly available at~\cite{thesisrepo}.

Writing every component from scratch is a deliberate choice. Production
systems are written in compiled languages and layered with
optimisations, batching, and low-level tricks that make them fast but
hard to follow. The goal here is the opposite: to expose clearly how a
transparent proof system is put together. A self-contained implementation
makes each step inspectable and keeps the distance between the
mathematics and the code as small as possible, at the natural cost of the
raw performance a production implementation would reach.

\section{Structure of the Thesis}
\label{sec:structure}

\Cref{ch:background} sets up the mathematical and cryptographic
background: finite fields, polynomials and the Number Theoretic
Transform, coding theory and Reed-Solomon codes, hash functions, proof
systems, and polynomial commitment schemes. \Cref{ch:iop} introduces the
interactive oracle proof model and the BCS transformation that compiles
it into a non-interactive argument using Merkle commitments.
\Cref{ch:fri} presents the FRI protocol, its folding operation, and its
soundness. \Cref{ch:pcs} turns FRI into a polynomial commitment scheme
through the quotient technique. \Cref{ch:implementation} describes the
implementation, its cryptographic components, its optimisations, and the
parameters and security level. \Cref{ch:conclusions} draws some concluding 
remarks.